\section{Proofs}\label{sec:proofs}

This is the section that is most like what a degree in maths would look like.

There are a few ways to attempt a proof.

\subsection{Proof by Contradiction} \label{sec:proof-contradiction}
A proof by contradiction works by assuming the opposite is true, and then applying logically correct steps until we reach a logical absurdity such as $1 = 0$. 
If we reached a logical absurdity via logically correct steps, then the original assumption must be wrong. Since the opposite of what we wanted to prove is wrong, 
the original statement was correct.

\begin{example}

    Prove by contradiction that 
    \[ (a - 2\sqrt{ab})(a + 2\sqrt{ab}) \geq - 4b^2\]
    for any non-negative real numbers $a$, $b$.
\end{example}

\begin{solution}
    First assume the opposite is true:
    \[(a - 2\sqrt{ab})(a + 2\sqrt{ab}) < - 4b^2 \]

    Then let us rearrange this until we reach a contradiction.

    \begin{align*}
        (a - 2\sqrt{ab})(a + 2\sqrt{ab}) &< - 4b^2 \\
        a^2 - 4ab &< - 4b^2 \tag{expanding}\\
        a^2 - 4ab + 4b^2  &< 0 \tag{rearranging}\\
        (a - 2b)^2 &< 0 \tag{factorising}
    \end{align*}

    This is a contradiction because the left-hand side is a square number, and a square number cannot be less than zero. Therefore our original assumption
    \[(a - 2\sqrt{ab})(a + 2\sqrt{ab}) < - 4b^2 \]
    was false, which means that the following must be true:

    \[ (a - 2\sqrt{ab})(a + 2\sqrt{ab}) \geq - 4b^2\]

    This completes the proof.
    \qed
\end{solution}


\subsection{Proof by Deduction}

Proof by deduction is the act of starting with a true statement, applying logically correct steps, and ending on a statement that must therefore be true.

It is like ``proof because it is true'' but made rigorous.

\begin{example}
    
    Starting with the statement 
    \[(a - 2b)^2 \geq 0\]
    Prove by deduction that 
    \[ (a - 2\sqrt{ab})(a + 2\sqrt{ab}) \geq - 4b^2\]
    for any non-negative real numbers $a$, $b$.
\end{example}
\begin{solution}
    We start with the true statement, and rearrange to deduce the desired final statement.
        \begin{align*}
        (a - 2b)^2 &\geq 0 \tag{starting with the hint} \\
        a^2 - 4ab + 4b^2  &\geq 0 \tag{expanding} \\
        a^2 - 4ab  &\geq -4b^2 \tag{rearranging} \\
        a^2 - \sqrt{4ab}^2 & \geq - 4b^2 \tag{using that $a$, $b$ are non-negative real numbers}\\
        (a - \sqrt{4ab})(a + \sqrt{4ab}) & \geq - 4b^2 \tag{difference of two squares}\\
        (a - 2\sqrt{ab})(a + 2\sqrt{ab})&\geq - 4b^2 \tag{using $\sqrt{xy} = \sqrt{x}\sqrt{y}$}
        \end{align*}
    Which is the statement we were trying to prove. This completes the proof.
    \qed
\end{solution}


\subsection{Proof by Exhaustion}

Proof by exhaustion is a method of proof that is essentially brute force; you try every case, and verify each case.

This can be used when there are a small number of cases to test.

Often, it is possible to use algebra to reduce what seems like a lot (even infinite) cases to only a few.

For example, if you need to prove that $n^2$ preserves even and odd for all integers, 
you can test all integers in two cases by splitting them into even $2m$ numbers and odd $2m-1$ numbers. 
In this way we reduce ``all integers'' to two cases and can prove it by exhaustion.

\begin{example}

    Prove by exhaustion that, for every integer $n$ whenever the left-hand side is defined, 
    \[ \sqrt{5-n^2} < 3 \]

\end{example}

\begin{solution}
    Since $\sqrt{5-n^2} $ is only defined when $5-n^2\geq0$, we only need to check the integers 
    $-2$, $-1$, $0$, $1$, $2$. 

    Since $(-n)^2 = n^2$, we only need to check $0$, $1$, and $2$,

    Starting with $n=2$,

    \begin{align*}
    \sqrt{5-2^2} &= \sqrt{5-4} \\
     &= \sqrt{1} \\
     &= 1
    \end{align*}
    And since $1<3$, the statement is true for $n=2$.

    For $n=1$,

    \begin{align*}
    \sqrt{5-1^2} &= \sqrt{5-1} \\
     &= \sqrt{4} \\
     &= 2
    \end{align*}
    And since $2<3$, the statement is true for $n=1$.

    For $n=0$,

    \begin{align*}
    \sqrt{5-0^2} &= \sqrt{5-0} \\
     &= \sqrt{5} 
    \end{align*}

    Now to check if $\sqrt{5}$ is less than $3$, we can observe that 

    \[ 5<9\]

    And therefore, taking square roots of both sides, we can deduce

    \[ \sqrt{5} < 3\]

    Which completes the proof by exhaustion, because we have checked every case in which 
    \[ \sqrt{5-n^2} < 3 \]

    is defined.
    \qed
\end{solution}

\begin{example}
    Using algebra, prove by exhaustion that $n^2-3$ is never a multiple of $9$ for any integer $n$.
\end{example}
\begin{solution}
    We can split integers into three cases: $3m$, $3m+1$, and $3m+2$ where $m$ is an integer.

    This accounts for every integer.

    Case 1: $n= 3m+2$

    Let's substitute and rearrange to see what we can deduce.

    \begin{align*}
        n^2-3 &= (3m+2)^2 - 3 \tag{using $n= 3m+2$}\\
         &= 9m^2 + 12m + 4 - 3 \tag{expanding}\\
         &=  9m^2 + 12m +1 \\
         &= 3(3m^2 + 4m) +1
    \end{align*}
    Looking at the final line, we see that when $n= 3m+2$, $n^2-3$ is one more than a multiple of $3$, and therefore never a multiple of $3$

    Since $9$ is a multiple of $3$, it is also never a multiple of $9$.

    Case 2: $n= 3m+1$

    Let's substitute and rearrange to see what we can deduce.

    \begin{align*}
        n^2-3 &= (3m+1)^2 - 3 \tag{using $n= 3m+1$}\\
         &= 9m^2 + 6m + 1 - 3 \tag{expanding}\\
         &=  9m^2 + 6m -2 \\
         &= 3(3m^2 + 2m) -2
    \end{align*}
    Looking at the final line, we see that when $n= 3m+1$, $n^2-3$ is two less than a multiple of $3$, and therefore never a multiple of $3$.

    Since $9$ is a multiple of $3$, it is also never a multiple of $9$.

    Case 3: $n=3m$

    Let's substitute and rearrange to see what we can deduce.

    \begin{align*}
        n^2-3 &= (3m)^2 - 3 \tag{using $n= 3m$}\\
         &= 9m^2 - 3 \tag{expanding}\\
    \end{align*}
    Looking at the final line, we see that when $n= 3m$, $n^2-3$ is three less than a multiple of $9$, and therefore never a multiple of $9$

    Since we have exhausted all cases, we can conclude that $n^2-3$ is never a multiple of $9$ for any integer $n$. This completes the proof.
    \qed
\end{solution}

